\documentclass[11pt,a4paper]{article}
\usepackage[utf8]{inputenc}
\usepackage[UTF8]{ctex} 
\usepackage{amsmath, amssymb, amsfonts, amsthm}
\usepackage{geometry}
\usepackage{float}
\usepackage{bm}
\usepackage{tcolorbox}
\usepackage[colorlinks=true, linkcolor=blue, unicode=true]{hyperref}

\geometry{left=2cm, right=2cm, top=2.5cm, bottom=2.5cm}

\title{论文笔记：K-均值聚类的分层初始化方法 \\ (Hierarchical Initialization for K-Means)}
\date{\today}

\newtheorem{property}{性质}

\begin{document}
\maketitle

\section{核心思想与背景}
该方法将聚类问题转化为\textbf{加权聚类问题}。通过类似金字塔结构的自底向上采样减少数据规模，在高层寻找初始中心，再自顶向下逐层精炼 。

\section{数学准备与关键性质}

\subsection{加权采样定义 (2D 示例)}
设第 $k$ 层有 4 个相邻点 $(i,j)$，其权值为 $w_{k,i,j}$。采样到第 $k+1$ 层后，产生新点 $P(P_x, P_y, P_w)$ ：
\begin{itemize}
    \item \textbf{坐标映射}：$P_x = \frac{i+1}{2}, P_y = \frac{j+1}{2}$ [cite: 519]。
    \item \textbf{权重计算}：$P_w = \frac{w_{k,i,j} + w_{k,i+1,j} + w_{k,i,j+1} + w_{k,i+1,j+1}}{4}$ 。
\end{itemize}

\subsection{核心理论性质}
\begin{property}[质心性质]
对于加权密度聚类，若使用平方误差准则，聚类中心即为其加权质心 。
\end{property}

\begin{property}[投影误差界限]
设 $X_{k+1}$ 为第 $k+1$ 层中心，$X_k$ 为第 $k$ 层对应中心，则满足 ：
\begin{equation}
    2 \cdot X_{k+1} - X_k \le 1
\end{equation}
该性质证明了高层中心投影回底层后，距离真实最优中心极近（误差 $\le 1$），保证了极快的收敛速度 。
\end{property}

\subsection{性质 2 的证明 (Proof of Property 2)}

\textbf{证明目标}：证明 $0 \le 2 \cdot X_{k+1,j} - X_{k,j} \le 1$。

\begin{proof}
1. \textbf{定义两层质心坐标}：
根据性质 1，第 $k+1$ 层某簇 $j$ 的质心坐标为：
\begin{equation}
    X_{k+1,j} = \frac{\sum_{m=1}^{n} W_{k+1,m} x_{k+1,m}}{\sum_{m=1}^{n} W_{k+1,m}} \quad \text{}
\end{equation}
对应的第 $k$ 层质心坐标为：
\begin{equation}
    X_{k,j} = \frac{\sum_{l=1}^{4n} w_{k,l} x_{k,l}}{\sum_{l=1}^{4n} w_{k,l}} \quad \text{}
\end{equation}

2. \textbf{统一坐标系}：
为了比较，将 $X_{k+1,j}$ 投影回第 $k$ 层（坐标乘以 2），计算差值：
\begin{equation}
    \Delta = 2 \cdot X_{k+1,j} - X_{k,j} = \frac{2 \cdot \sum_{m=1}^{n} W_{k+1,m} x_{k+1,m}}{\sum_{m=1}^{n} W_{k+1,m}} - \frac{\sum_{l=1}^{4n} w_{k,l} x_{k,l}}{\sum_{l=1}^{4n} w_{k,l}} \quad \text{}
\end{equation}

3. \textbf{权重代换}：
设第 $k$ 层所有相关点的权重总和为 $T = \sum w_{k,l}$，第 $k+1$ 层为 $S = \sum W_{k+1,m}$。
由于 $W_{k+1,m} = \frac{w_{k,i} + w_{k,i+1} + w_{k,j} + w_{k,j+1}}{4}$，易证 $T = 4S$ 。
由此可将差值公式简化为：
\begin{equation}
    \Delta = \frac{8 \cdot \sum_{m=1}^{n} W_{k+1,m} x_{k+1,m} - \sum_{l=1}^{4n} w_{k,l} x_{k,l}}{T} \quad \text{}
\end{equation}

4. \textbf{单点分析}：
考虑第 $k+1$ 层的一个点及其对应的第 $k$ 层四个点（坐标分别为 $i, i+1$）。
第 $k+1$ 层项贡献为：$8 \cdot W_{k+1,m} \cdot \frac{i+1}{2} = (w_{k,i,j} + w_{k,i+1,j} + w_{k,i,j+1} + w_{k,i+1,j+1})(i+1)$ 。
第 $k$ 层项贡献为：$w_{k,i,j} \cdot i + w_{k,i+1,j}(i+1) + w_{k,i,j+1} \cdot i + w_{k,i+1,j+1}(i+1)$ 。
两者相减得：
\begin{equation}
    \text{分子单项差} = w_{k,i,j} + w_{k,i,j+1} \quad \text{}
\end{equation}

5. \textbf{结论}：
将所有点累加，分子变为部分原始权重的和，而分母 $T$ 是所有权重的和。
由于权重 $w \ge 0$，显然有：
\begin{equation}
    0 \le \frac{\sum (w_{k,i,j} + w_{k,i,j+1})}{T} \le 1 \quad \text{}
\end{equation}
证毕。
\end{proof}

\section{算法详细流程}

\begin{tcolorbox}[title=分层初始化算法步骤]
\begin{description}
    \item[Step 0: 预处理] 将数据坐标线性映射为 $1 \dots 2^N$ 的整数。
    \item[Step 1: 自底向上采样] 重复执行平均采样，直到数据量减小至 $\ge 20 \cdot C_{num}$ (即簇数的20倍) 。
    \item[Step 2: 顶层初始化] \textbf{排序策略}：选取权值最大的 $C_{num}$ 个点作为初始中心 。此举利用了采样过程的“平滑效应”来抑制噪声。
    \item[Step 3: 自顶向下精炼] 将 $k+1$ 层中心坐标翻倍投影至 $k$ 层，作为初始值进行迭代，直至回到原始层 。
    \item[Step 4: 后处理] 执行逆变换回到原始坐标空间 。
\end{description}
\end{tcolorbox}

\section{高维情况处理 ($D > 2$)}
针对高维数据的稀疏性，算法进行了以下优化 ：
\begin{itemize}
    \item \textbf{存储优化}：不使用全空间网格，仅使用 $N \times D$ 数组存储非零点 。
    \item \textbf{权重合并}：采样后，若不同原始点映射到了相同的整数坐标，则将其权值累加 。
    \item \textbf{权值衰减}：高维采样时，单点权值需除以 $2^D$ 。
\end{itemize}

\section{实验结论与性能对比}
\begin{itemize}
    \item \textbf{时间效率}：在 Iris 和 Glass 数据集上，由于初始化极佳，每一层迭代通常少于 10 次 。总运行时间显著低于 KA 方法和 Fayyad 的随机采样法 。
    \item \textbf{抗噪性}：采样过程类似于低通滤波器，能有效识别数据“峰值”并抑制分布均匀的噪声 。
    \item \textbf{局限性}：实验表明，在维度较低时（$<10$ 维）优于 Fayyad 法；但在极高维度下，性能可能略逊于 Fayyad 法 。
\end{itemize}

\end{document}
